\section{Source Section II.C: Equilibrium Conditions and Wall Terms}
\label{sec:equilibrium-conditions}

\sourceeqrange{Eqs. (2.21)--(2.34)}

\subsection{Purpose and Scope}

This section continues the source order after the gradient entropy and energy
construction.  It explains the equilibrium variational branch, the
one-dimensional interface first integral, the surface-tension change of
variables, and the natural wall density boundary condition.  The derivation is
written as a companion guide; the source paper remains the required reference
for the original presentation.

The scope is deliberately bounded.  The reversible stress tensor and
hydrodynamic equations are not derived here, and Appendix-B scaling is not used
in this section.

\subsection{Typed Objects and Assumptions}

Use the Section II.B fields \(n\), \(\hat e\), \(T\), \(M\), and \(\hat\mu\).
For equilibrium, the admissible states are time-independent fields on a
container \(\Omega\), with \(n>0\), \(T>0\), and enough smoothness for the
Euler--Lagrange and boundary calculations below.

Additional objects introduced in Eqs.~(2.21)--(2.34) are:

\begin{description}
\item[\(N=\int_\Omega n\dd x\)] scalar functional for the total particle
  number constraint.
\item[\(W\)] auxiliary constrained entropy functional.
\item[\(\beta\)] scalar energy-constraint multiplier.
\item[\(\lambda_N\)] scalar particle-number multiplier in companion notation.
\item[\(F\)] bulk Helmholtz free-energy functional after \(T\) is homogeneous.
\item[\(f(n,T)\)] local Helmholtz free-energy density without the gradient
  term.
\item[\(g_{\mathrm{gp}}(n,T)\)] grand-potential density for a planar interface.
\item[\(\gamma(T)\)] surface tension for the planar branch.
\item[\(S_s,E_s\)] boundary functionals for wall entropy and wall energy.
\item[\(\sigma_s,e_s\)] surface entropy and surface energy densities.
\item[\(f_s=e_s-T\sigma_s\)] wall Helmholtz free-energy density.
\item[\(a_s,b_s,n_c\)] wall parameters in the quadratic wall model.
\end{description}

The source symbol for the particle-number multiplier is close to the normal
symbol.  The companion writes this multiplier as \(\lambda_N\) to avoid
confusing it with the outward unit normal \(\nu\).

\paragraph{Equilibrium and wall status.}
The Section II.C objects are equilibrium objects unless the text says
otherwise.  This ledger records which statements are source-stated model choices
and which are derived within the companion.
\begin{description}[leftmargin=2.8cm,style=nextline]
\item[\(W,N,\beta,\lambda_N\)] source-stated constrained-equilibrium setup;
no dynamics or transport law is introduced here.
\item[\(T=1/(k_B\beta)\)] Derived stationarity condition for energy variations
on the homogeneous equilibrium branch.
\item[\(F\), \(g_{\mathrm{gp}}\), first integral] Derived free-energy and
planar-interface relations requiring coexistence endpoints and a
one-dimensional monotone branch.
\item[\(S_s,E_s,\sigma_s,e_s\)] source-stated wall entropy/energy model
objects; they are not supplied by the bulk variation.
\item[\(f_s\), \(F_{\mathrm{tot}}\), Eq.~\eqref{eq:iic-natural-wall-condition}]
Derived equilibrium wall free-energy branch requiring free boundary density
variation or an equivalent natural branch.
\item[Dynamic wall transfer] Heat, entropy, and total-energy transfer remain
unresolved later-source data; they are not closed by the equilibrium wall
condition.
\end{description}


\subsection{Section II.B Boundary Term Reused Here}

Section II.B gave the entropy variation boundary contribution
\begin{equation}
  -\int_{\partial\Omega}
  \frac{M}{T}(\nu\cdot\grad n)\delta n\,\dd a .
  \label{eq:iic-iib-boundary}
\end{equation}
At homogeneous equilibrium temperature, the corresponding bulk free-energy
boundary contribution is
\begin{equation}
  \int_{\partial\Omega}
  M(\nu\cdot\grad n)\delta n\,\dd a .
  \label{eq:iic-free-energy-boundary}
\end{equation}
This sign is the one used below for the wall boundary condition.  The passage
from \eqref{eq:iic-iib-boundary} to \eqref{eq:iic-free-energy-boundary} is not
a global entropy-balance statement; it is only the equilibrium free-energy
variation branch.

\subsection{Constrained Equilibrium Functional}

The equilibrium problem first fixes total particle number \(N\) and bulk
internal energy \(E_b\).  The two constraints play different roles in the
step-by-step calculation:
\begin{center}
\begin{tabular}{@{}lll@{}}
\toprule
Constraint & Multiplier & Variation it controls \\
\midrule
particle number \(N\) & \(\lambda_N\) & density variation \(\delta n\) \\
bulk energy \(E_b\) & \(\beta\) & energy variation \(\delta e\) \\
\bottomrule
\end{tabular}
\end{center}
In companion notation, the constrained entropy
functional is
\begin{equation}
  W
  =
  \frac{S_b}{k_B}
  +
  \lambda_N N
  -
  \beta E_b .\sourceeqbadge{2.21}
  \label{eq:iic-W}
\end{equation}
Varying \(e\) at fixed \(n\) gives the energy part of the first variation,
\begin{equation}
  \delta_e W
  =
  \int_\Omega
  \left(\frac{1}{k_B T}-\beta\right)\delta e\,\dd x .
  \label{eq:iic-W-energy-variation}
\end{equation}
Here \(1/T=n(\partial s/\partial e)_n\) is the local Section II.B
entropy-density coefficient.  Since \(\delta e\) is arbitrary in the bulk,
stationarity gives a spatially constant temperature,
\begin{equation}
  T=\frac{1}{k_B\beta}.\sourceeqbadge{2.22}
  \label{eq:iic-homogeneous-temperature}
\end{equation}

Once \(T\) is homogeneous, \(F=E_b-TS_b\) is the bulk Helmholtz free-energy
functional.  Substituting \(\beta=1/(k_B T)\) into \eqref{eq:iic-W} gives the
source rewriting
\begin{align}
  W
  &=\frac{S_b}{k_B}+\lambda_N N-\frac{1}{k_B T}E_b \\
  &=\frac{T S_b-E_b}{k_B T}+\lambda_N N \\
  &=\frac{-F+k_B T\lambda_N N}{k_B T}.\sourceeqbadge{2.23}
  \label{eq:iic-W-helmholtz-conversion}
\end{align}
Thus the stationarity of \(W\) is equivalent on this branch to stationarity of
\(F-k_B T\lambda_N N\).  Equivalently, for variations at homogeneous
\(T\),
\[
  -k_BT\,\delta W
  =\delta(E_b-TS_b)-k_BT\lambda_N\delta N
  =\delta F-k_BT\lambda_N\delta N,
\]
where \(\delta T=0\) on this homogeneous-temperature equilibrium branch.  The
bulk Helmholtz functional follows from \(F=E_b-TS_b\) and the Section II.B
functionals:
\begin{align}
  F[n]
  &=\int_\Omega\left(e+\frac{1}{2}K|\grad n|^2\right)\dd x
    -T\int_\Omega\left(ns-\frac{1}{2}C|\grad n|^2\right)\dd x \\
  &=\int_\Omega
  \left(
    f(n,T)+\frac{1}{2}(K+TC)|\grad n|^2
  \right)\dd x \\
  &=\int_\Omega
  \left(
    f(n,T)+\frac{1}{2}M(n,T)|\grad n|^2
  \right)\dd x .\sourceeqbadge{2.24}
  \label{eq:iic-F}
\end{align}
Then maximizing \(W\) with respect to \(n\) is equivalent to stationarity of
\(F-k_B T\lambda_N N\).  The equivalence follows from multiplying the
stationarity condition for \(W\) by the positive constant \(-k_BT\) after the
energy variation has fixed a homogeneous temperature.  Thus the density part is
an equilibrium free-energy variation, not a dynamic entropy-production law.
The density variation is
\begin{equation}
  \delta F
  =
  \int_\Omega \hat\mu\,\delta n\,\dd x
  +
  \int_{\partial\Omega}M(\nu\cdot\grad n)\delta n\,\dd a ,
  \label{eq:iic-delta-F}
\end{equation}
where the bulk coefficient is the generalized chemical potential from
Section II.B at homogeneous \(T\).  Including the particle-number constraint,
\begin{align}
  \delta\left(F-k_B T\lambda_N N\right)
  &=\int_\Omega (\hat\mu-k_B T\lambda_N)\delta n\,\dd x \\
  &\quad+
  \int_{\partial\Omega}M(\nu\cdot\grad n)\delta n\,\dd a .
  \label{eq:iic-density-stationarity-expanded}
\end{align}
The phrase "controlled separately" has three concrete meanings used later in
this companion:
\begin{center}
\begin{tabular}{@{}lll@{}}
\toprule
Boundary branch & Boundary variation & Consequence \\
\midrule
Fixed trace & \(\delta n=0\) on \(\partial\Omega\) & no natural wall condition follows \\
Natural wall & \(\delta n\) free & boundary coefficient must vanish \\
Wall functional & \(\delta n\) free with \(f_s\) present & natural condition includes \((\partial f_s/\partial n)_T\) \\
\bottomrule
\end{tabular}
\end{center}
If the boundary term is controlled by one of these branches, the bulk
Euler--Lagrange condition is
\begin{equation}
  \hat\mu=k_B T\lambda_N=\mathrm{constant}.\sourceeqbadge{2.25}
  \label{eq:iic-muhat-constant}
\end{equation}
This is a local stationarity condition.  It is not a hydrodynamic balance law.
On this equilibrium branch, the generalized chemical potential is constant in
the bulk.

\subsection{One-Dimensional Interface Branch}

For a planar interface, write \(n=n(x)\) and assume a homogeneous temperature.
The companion uses
\begin{equation}
  g_{\mathrm{gp}}(n,T)
  =
  f(n,T)-\mu_{\mathrm{cx}}n+p_{\mathrm{cx}} .\sourceeqbadge{2.26}
  \label{eq:iic-grand-potential-density}
\end{equation}
for the source grand-potential density.  The subscript is added only to avoid
confusion with the gravitational acceleration used later.  At coexistence,
\(p_{\mathrm{cx}}=\mu_{\mathrm{cx}}n-f\) for each bulk endpoint, so
\(g_{\mathrm{gp}}(n_g,T)=g_{\mathrm{gp}}(n_l,T)=0\).

The one-dimensional functional density for the planar branch is
\[
  L(n,n_x)
  =
  g_{\mathrm{gp}}(n,T)
  +
  \frac{1}{2}M(n,T)n_x^2 .
\]
Because \(L\) has no explicit \(x\)-dependence, the Beltrami identity gives a
first integral.  One common sign convention writes \(L-n_x L_{n_x}=\text{constant}\);
the companion uses the negative of that constant so that the far-field value is
zero in the same orientation as the source formula.  The derivative with
respect to the interface slope is
\[
  \frac{\partial L}{\partial n_x}=M n_x,
\]
so
\begin{equation}
  n_x\frac{\partial L}{\partial n_x}-L
  =
  M n_x^2-\left(g_{\mathrm{gp}}+\frac{1}{2}M n_x^2\right)
  =
  \frac{1}{2}M n_x^2-g_{\mathrm{gp}}
  =
  \mathrm{constant}.
  \label{eq:iic-first-integral-raw}
\end{equation}
In the coexistence far fields, \(n_x\to0\) and
\(g_{\mathrm{gp}}\to0\), so the constant is zero.  Therefore the planar
interface branch satisfies
\begin{equation}
  g_{\mathrm{gp}}(n,T)
  =
  \frac{1}{2}M(n,T)\left(\frac{\dd n}{\dd x}\right)^2 .
  \label{eq:iic-first-integral}
\end{equation}
This identity is the source's first-integral statement following Eq.~(2.26).  It belongs to the equilibrium one-dimensional branch only.

\subsection{Surface-Tension Identity}

The surface tension is an excess free energy per unit interfacial area.  Along
the planar branch,
\[
  \gamma
  =
  \int_{-\infty}^{\infty}
  \left[
    g_{\mathrm{gp}}(n,T)
    +
    \frac{1}{2}M(n,T)n_x^2
  \right]\dd x .
\]
Using \eqref{eq:iic-first-integral}, this becomes
\begin{equation}
  \gamma
  =
  \int_{-\infty}^{\infty}M(n,T)n_x^2\,\dd x .
  \label{eq:iic-gamma-x}
\end{equation}
For a monotone interface, \(\dd n=n_x\dd x\), and
\eqref{eq:iic-first-integral} gives
\[
  |n_x|=\sqrt{\frac{2g_{\mathrm{gp}}}{M}} .
\]
Choosing the orientation from gas density \(n_g\) to liquid density \(n_l\)
with \(n_x>0\), the integrand transforms as
\begin{equation}
  M n_x^2\,\dd x
  =M n_x\,\dd n
  =\sqrt{2M g_{\mathrm{gp}}}\,\dd n .
  \label{eq:iic-surface-tension-change-variable-step}
\end{equation}
The change of variables gives
\begin{equation}
  \gamma(T)
  =
  \int_{n_g}^{n_l}
  \left[2M(n,T)g_{\mathrm{gp}}(n,T)\right]^{1/2}\dd n .\sourceeqbadge{2.27}
  \label{eq:iic-gamma-n}
\end{equation}
This derivation uses the equilibrium first integral and monotonicity.  It is
not a statement about a dynamic diffuse interface.

\subsection{Surface Entropy and Surface Energy}

The wall branch adds boundary functionals
\begin{equation}
  S_s[n_{\partial\Omega}]
  =
  \int_{\partial\Omega}\sigma_s(n)\dd a,
  \qquad
  E_s[n_{\partial\Omega}]
  =
  \int_{\partial\Omega}e_s(n)\dd a .\sourceeqbadge{2.28}
  \label{eq:iic-surface-functionals}
\end{equation}
The densities \(\sigma_s\) and \(e_s\) are functions of the boundary trace of
\(n\).  They are source-stated wall objects, not consequences of the bulk
calculation alone.  With the source's surface-entropy convention, the total
constrained wall functional is
\begin{equation}
  W_{\mathrm{tot}}
  =
  W+
  \int_{\partial\Omega}\left(\sigma_s(n)-\beta e_s(n)\right)\dd a .\sourceeqbadge{2.29}
  \label{eq:iic-Wtot}
\end{equation}
Here the surface term is written in the same scaled convention as source
Eq.~(2.29); the companion does not silently replace it by a differently scaled
surface entropy.

At homogeneous equilibrium temperature, define the wall Helmholtz free-energy
density
\begin{equation}
  f_s(n,T)=e_s(n)-T\sigma_s(n).\sourceeqbadge{2.31}
  \label{eq:iic-fs}
\end{equation}
The total equilibrium free-energy functional for the wall branch is
\begin{equation}
  F_{\mathrm{tot}}
  =
  F
  +
  \int_{\partial\Omega}f_s(n,T)\dd a .\sourceeqbadge{2.32}
  \label{eq:iic-Ftot}
\end{equation}

\subsection{Natural Wall Density Boundary Condition}

Vary \(F_{\mathrm{tot}}\) with respect to the boundary trace of \(n\).  The
bulk gradient term contributes \eqref{eq:iic-free-energy-boundary}.  The wall
functional contributes
\[
  \int_{\partial\Omega}
  \left(\frac{\partial f_s}{\partial n}\right)_T
  \delta n\,\dd a .
\]
The boundary coefficient is assembled from exactly two sources:
\[
  \underbrace{M\nu\cdot\grad n}_{\text{bulk gradient term}}
  +
  \underbrace{(\partial f_s/\partial n)_T}_{\text{wall free energy}}.
\]
Thus the boundary part of the first variation is
\begin{equation}
  \delta F_{\mathrm{tot}}\big|_{\partial\Omega}
  =
  \int_{\partial\Omega}
  \left[
    M\nu\cdot\grad n
    +
    \left(\frac{\partial f_s}{\partial n}\right)_T
  \right]\delta n\,\dd a .
  \label{eq:iic-wall-boundary-variation-expanded}
\end{equation}
For free boundary variations, the coefficient of \(\delta n\) on
\(\partial\Omega\) must vanish:
\begin{equation}
  M\nu\cdot\grad n
  +
  \left(\frac{\partial f_s}{\partial n}\right)_T
  =
  0 .\sourceeqbadge{2.30}
  \label{eq:iic-natural-wall-condition}
\end{equation}
This is the natural wall density boundary condition for the equilibrium wall
functional.  If \(\delta n\) is fixed on the boundary instead, the same
condition is not forced by variation.

For the quadratic wall model
\begin{equation}
  f_s(n,T)
  =
  -a_s(n-n_c)+\frac{1}{2}b_s(n-n_c)^2 ,
  \sourceeqbadge{2.33}
  \label{eq:iic-quadratic-fs}
\end{equation}
the derivative is a one-line ordinary derivative with the wall trace of \(n\) as
its scalar argument:
\begin{align}
  \frac{\partial}{\partial n}\{-a_s(n-n_c)\}&=-a_s, \\
  \frac{\partial}{\partial n}\left\{\frac{1}{2}b_s(n-n_c)^2\right\}&=b_s(n-n_c).
  \label{eq:iic-quadratic-fs-derivative-ledger}
\end{align}
Therefore
\[
  \left(\frac{\partial f_s}{\partial n}\right)_T
  =
  -a_s+b_s(n-n_c).
\]
Substitution into \eqref{eq:iic-natural-wall-condition} gives
\begin{equation}
  M\nu\cdot\grad n-a_s+b_s(n-n_c)=0 .\sourceeqbadge{2.34}
  \label{eq:iic-quadratic-wall-condition}
\end{equation}

\subsection{Wall Branch Versus Global Transfer}

The wall condition above is an equilibrium boundary condition.  It does not by
itself close the later dynamic global entropy or total-energy balance.  That
later closure also needs heat-transfer data, surface-energy time variation, and
the hydrodynamic balance laws.  The companion therefore keeps three statements
separate:

\begin{itemize}
\item the local Section II.B variation boundary term;
\item the equilibrium wall density condition derived from \(F_{\mathrm{tot}}\);
\item the later dynamic global entropy and energy transfer branch.
\end{itemize}

\subsection{Verification Checklist}

The finite checks for this section cover:

\begin{enumerate}
\item the one-dimensional integration-by-parts identity for the equilibrium
  gradient term;
\item the boundary residual
  \(M\nu\cdot\grad n+(\partial f_s/\partial n)_T\) and its cancellation under
  the natural wall condition;
\item the derivative of the quadratic wall free-energy density;
\item the algebraic consistency of the planar surface-tension integrand under
  \(g_{\mathrm{gp}}=Mn_x^2/2\);
\item the expected nonzero residual when the wall contribution is omitted while
  the boundary variation is free;
\item regression text that global entropy transfer remains unresolved here.
\end{enumerate}

These checks verify finite variational bookkeeping.  They do not prove the wall
model, PDE well-posedness, hydrodynamic equations, or simulation results.

The executable mirrors for the Section II.C calculations are:
\begin{center}
\small
\begin{tabular}{@{}l@{}}
\texttt{scripts/python/section\_iic\_equilibrium\_wall\_sympy.py}\\
\texttt{scripts/wolfram/section\_iic\_equilibrium\_wall\_checks.wls}\\
\texttt{tests/test\_section\_iic\_equilibrium\_wall.py}
\end{tabular}
\end{center}
They check the equilibrium integration-by-parts residual, the planar
first-integral algebra, the wall free-energy derivative, and the expected
nonzero residual when the wall contribution is removed.  These are finite
checks of the displayed companion steps, not a closure of the later dynamic
global-transfer problem.


\subsection{Source-Procedure Closure for Section II.C}

The completion audit treats the Section II.C equilibrium branch as covered when
the constrained-temperature step, planar first integral, wall functional, and
natural boundary condition are each explicit.  This closure is equilibrium-only;
it does not supply dynamic heat, entropy, or surface-energy transfer terms.

\begin{center}
\begin{tabular}{@{}p{0.18\linewidth}p{0.30\linewidth}p{0.38\linewidth}@{}}
\toprule
Source item & Procedure type & Companion coverage \\
\midrule
\eqsrc{2.21}--\eqsrc{2.25} & constrained entropy and free-energy stationarity & energy variation fixes homogeneous \(T\); Eq.~(2.23) rewrites \(W\); Eq.~(2.24) defines \(F\); Eq.~(2.25) gives constant \(\hat\mu\) \\
\eqsrc{2.26}--\eqsrc{2.27} & planar first integral and surface tension & the grand-potential density is Eq.~(2.26); the Beltrami first integral is source text before Eq.~(2.27); the density integral is Eq.~(2.27) \\
\eqsrc{2.28}--\eqsrc{2.32} & wall entropy/energy, wall constraint, and total free energy & Eq.~(2.29) is the constrained wall functional; Eq.~(2.30) is the natural boundary condition; Eq.~(2.32) is \(F_{\mathrm{tot}}\) \\
\eqsrc{2.33}--\eqsrc{2.34} & quadratic wall branch & the source wall model is differentiated and substituted into Eq.~(2.30) \\
\bottomrule
\end{tabular}
\end{center}

The source rows for total wall entropy and energy are therefore covered as
static equilibrium wall functionals.  They remain separate from the later
dynamic global entropy and total-energy balance.

\subsection*{Tagged verification complement}

\OnukiLinkIIC{The tagged Section II.C verification complement} gives the
equilibrium hand route and finite boundary checks, with the dynamic wall-transfer
gap retained explicitly.

\subsection{Open Issues Carried Forward}

\begin{itemize}
\item The reversible stress tensor is deferred to the hydrodynamic-equation
  section.
\item The hydrodynamic balance laws are not derived in this equilibrium
  section.
\item Global entropy transfer and total-energy transfer are not closed by the
  wall density condition alone.
\item Appendix-B B3 scaling is not involved in this section.
\end{itemize}
